% TOP_PROBLEM: {"id": 3006, "title": "Kirby Problem 4.130", "category_id": 7, "category": {"id": 7, "name": "topology", "display_name": "Topology", "description": "Properties preserved under continuous deformations.", "slug": "topology", "order_index": 7, "created_at": "2026-07-31T15:26:25.671Z"}, "status": "open"}
% FIRST_POSTED: null
% ENTRY_KIND: statement_only
% Imported from: batch14.tex; UnsolvedMath ID 3006
\documentclass[11pt]{article}
\usepackage[margin=25mm]{geometry}
\usepackage{fontspec}
\setmainfont{Latin Modern Roman}
\usepackage{amsmath,amssymb,mathtools}
\usepackage{unicode-math}
\setmathfont{Latin Modern Math}
\usepackage{xurl,hyperref}
\hypersetup{colorlinks=true,urlcolor=blue}
\setcounter{secnumdepth}{0}
\setlength{\parindent}{0pt}
\setlength{\parskip}{5pt}
\setlength{\emergencystretch}{3em}
\newcommand{\sref}[2]{\hyperlink{source:#1:#2}{[S#2]}}
\newcommand{\cataloguescope}[1]{\paragraph{Recorded target.}#1}
\begin{document}
% UnsolvedMath ID 3006; source catalogue code KP-4.130
\setcounter{section}{3005}
\section{Kirby Problem 4.130}\label{3006}
{\small\sffamily UnsolvedMath ID 3006 · Topology}
\subsection{Definitions and mathematical statement}

\paragraph{Shared notation.}$\mathbb N=\{1,2,\ldots\}$, and $\mathbb Z,\mathbb Q,\mathbb R,\mathbb C$ denote the integers, rationals, reals and complex numbers. Any other conventions are specified locally.


Let $X$ be a K3 surface with an elliptic fibration and let $T$ be a regular torus fiber, with the product framing $T\times D^2$. For a smooth knot $K\subset S^3$, let $E(K)=S^3\setminus\operatorname{int}\nu K$ be its exterior, with knot meridian $\mu_K$ and preferred (zero-framed) longitude $\lambda_K$ on its boundary. With fiber coordinate circles $a,b$ and normal meridian $\mu_T$ on $\partial(T\times D^2)$, longitudinal surgery uses the product gluing that identifies $\mu_K,\lambda_K$ with $a,b$ and the auxiliary circle $s$ in $S^1\times E(K)$ with $\mu_T$, with signs chosen for the oriented boundary gluing. Denote the resulting closed smooth manifold by
\[
 X_K^\lambda=(X\setminus\operatorname{int}(T\times D^2))\cup_\phi(S^1\times E(K)).
\]
This is the zero-framed knot surgery convention in the original construction. \sref{3006}{3}
A closed oriented smooth four-manifold is reducible if it admits a smooth connected-sum decomposition $M_1\# M_2$ with neither $M_i$ a homotopy four-sphere, meaning neither is homotopy equivalent to $S^4$. \sref{3006}{4} It is completely decomposable if it is diffeomorphic to a connected sum of copies of $\mathbb{CP}^2$ and $\overline{\mathbb{CP}}^{2}$.
\paragraph{Statement recorded in the supplied source.}
For every knot $K\subset S^3$, is the longitudinally surgered K3 manifold $X_K^\lambda$ reducible? Is it always completely decomposable? These are two separate conclusions about this gluing convention, which differs from ordinary Fintushel--Stern knot surgery. \sref{3006}{2}
\subsection{Short English statement}
Does longitudinal knot surgery on a regular fiber of a K3 surface always produce a connected sum with neither summand a homotopy four-sphere, and does it always split entirely into projective-plane summands?
\subsection{Sources}
\begin{description}
\item[\hypertarget{source:3006:1}{S1}] ULAM.AI and the UnsolvedMath Contributors, \emph{UnsolvedMath: A Curated Collection of Open Mathematics Problems}, record 3006 (KP-4.130). CC BY 4.0. \url{https://huggingface.co/datasets/ulamai/UnsolvedMath}.
\item[\hypertarget{source:3006:2}{S2}] R. İnanç Baykur, Robion C. Kirby and Daniel Ruberman (eds.), \emph{K3: A New Problem List in Low-Dimensional Topology}, Mathematical Surveys and Monographs 295, American Mathematical Society (2026), Problem 4.130, author preliminary version. \url{https://math.berkeley.edu/sites/default/files/surv-295-ruberman-watermarked-author-pdf.pdf}.
\item[\hypertarget{source:3006:3}{S3}] Clifford Henry Taubes, \emph{Some 4-manifold geometry from hyperbolic knots in $S^3$}, arXiv:\allowbreak1602.01687 (2016), §2(a), equation (2.1). \url{https://arxiv.org/abs/1602.01687}.
\item[\hypertarget{source:3006:4}{S4}] M. J. D. Hamilton and D. Kotschick, \emph{Minimality and irreducibility of symplectic four-manifolds}, International Mathematics Research Notices (2006), Article ID 35032, §1.2, p.~2 (arXiv:math/0509204v2). \url{https://arxiv.org/abs/math/0509204v2}.
\end{description}
\label{end:3006}
\end{document}
